% =============================================================
%  Chapter 5 — Synthesis: A Hierarchy of Complexity in
%              Continual Hamiltonian Learning
%  Rewritten as a trace-based synthesis chapter.
%  Primary structure: trace → theorem as summary → comparison table.
%  Every claim is grounded in equation/algorithm references
%  from Chapters 3 and 4.
% =============================================================

\chapter{Synthesis: A Hierarchy of Complexity in Continual Hamiltonian Learning}
\label{chap:hierarchy}

\providecommand{\Ham}{\mathcal{H}}
\providecommand{\Potential}{\mathcal{R}}
\providecommand{\PhaseState}{z}
\providecommand{\Pos}{q}
\providecommand{\Mom}{p}
\providecommand{\Mass}{\mathrm{M}}
\providecommand{\J}{\mathbb{J}}

% =======================================================================
% \section{Purpose of This Chapter}
% \label{sec:ch5:purpose}
% =======================================================================

Chapters~\ref{chap:grlsnam} and~\ref{chap:material} should not be read
as two independent navigation systems.
They share the same update skeleton; the primary dynamical difference
appears as one additional term in the $p$-update.
This chapter makes that precise.

The argument proceeds in three steps.
First, the shared skeleton: both stages have an identical $(p,q)$-update
structure; what changes between them is a single additive term in the
$p$-update.
Second, the trace: that single term forces, by necessity, five downstream
consequences: a new force channel, a new coefficient, a new objective
term, a new gradient pathway, and a new update loop.
Third, the theorem: the enrichment principle is the compact statement
of why one additive line costs five downstream consequences, and
Corollary~\ref{cor:ch5:concrete} verifies it by pointing to specific
equations in Chapters~\ref{chap:grlsnam} and~\ref{chap:material}.

% =======================================================================
\section{The Shared Skeleton: One p-Update, Two Instances}
\label{sec:ch5:skeleton}
% =======================================================================

Both stages share an identical phase-space update structure.
Written side by side:

\begin{equation}
\renewcommand{\arraystretch}{1.6}
\begin{array}{ll}
\textbf{Stage~1} & \textbf{Stage~2} \\
\hline
\Mom_{t+1} = \Mom_t
  - \tau\underbrace{\nabla_\Pos H_{\mathrm{geom}}}_{\text{geom. force}}
  - \tau R_\theta
  + \tau\mathbf{B}u_t
&
\Mom_{t+1} = \Mom_t
  - \tau\underbrace{\nabla_\Pos H_{\mathrm{geom}}}_{\text{geom. force}}
  - \tau\underbrace{\nabla_\Pos H_{\mathrm{mat}}}_{\mathbf{\text{new}}}
  - \tau R_\theta
  + \tau\mathbf{B}u_t \\[4pt]
\vec{q}_{t+1} = \vec{q}_t + \tau\Mass^{-1}\Mom_{t+1}
&
\vec{q}_{t+1} = \vec{q}_t + \tau\Mass^{-1}\Mom_{t+1}
\end{array}
\label{eq:ch5:side_by_side}
\end{equation}

The $q$-update is identical.
The $R_\theta$ port is identical.
The geometry force $\nabla_\Pos H_{\mathrm{geom}}$ is identical and
carried forward unchanged.
The \emph{only structural difference} is the single additive term
$-\tau\nabla_\Pos H_{\mathrm{mat}}$ in the Stage~2 $p$-update.

\begin{definition}[Continual Hamiltonian Navigation Instance]
\label{def:ch5:instance}
A \emph{continual Hamiltonian navigation instance} is a tuple
$\mathfrak{S} = (H_\theta, R_\theta, f_\xi, \mathcal{L}, \mathcal{U})$
where $H_\theta : T^*\mathcal{Q}\to\mathbb{R}$ is the stored energy,
$R_\theta$ the dissipative port, $f_\xi$ the meta-navigator,
$\mathcal{L}$ the objective, and $\mathcal{U}$ the update hierarchy.
The policy is the time-$\Delta t$ flow of the port-Hamiltonian update
induced by $H_\theta$, $R_\theta$, and control input $u$
(Definition~\ref{def:passivity}), and the
shared template is:
\begin{equation}
  \PhaseState_{t+1} = \Phi_{\vartheta_t}(\PhaseState_t, y_t),
  \qquad
  \vartheta_t = f_\xi(y_t, \vec{g}).
  \label{eq:ch5:template}
\end{equation}
\end{definition}

Table~\ref{tab:ch5:instances} shows the two stages.

\begin{table}[t]
\centering
\caption{%
  \textbf{Two instantiations of the shared skeleton.}
  The $(p,q)$-update structure \eqref{eq:ch5:side_by_side} is
  identical; all variation is in the contents of
  $H_\theta$, $\mathcal{L}$, $\mathcal{U}$.
  Equation references point to the specific machinery in
  Chapters~\ref{chap:grlsnam}--\ref{chap:material}.
}
\label{tab:ch5:instances}
\small
\setlength{\tabcolsep}{4pt}
\renewcommand{\arraystretch}{1.35}
\resizebox{\textwidth}{!}{
\begin{tabular}{lll}
\toprule
\textbf{Component} &
\textbf{Stage~1 (Ch.~\ref{chap:grlsnam})} &
\textbf{Stage~2 (Ch.~\ref{chap:material})} \\
\midrule
$\nabla_\Pos H_\theta$ in $p$-update
  & $\nabla_\Pos H_{\mathrm{geom}}$ only
    \eqref{eq:ch3:ham}
  & $\nabla_\Pos H_{\mathrm{geom}} + \nabla_\Pos H_{\mathrm{mat}}$
    \eqref{eq:ch4:ham} \\
$R_\theta$ (dissipative port)
  & $\gamma\Mass^{-1}\Mom$ (implicit)
  & $\gamma\Mass^{-1}\Mom$ (explicit separate port)
    \eqref{eq:ch4:R} \\
Context $y_t$
  & $\{d_j,\vec{n}_j\}$ geometry only
    \eqref{eq:ch3:context}
  & $\{d_j,\vec{n}_j\}\cup\{\tilde{r},\phi,P_t\}$
    \eqref{eq:ch4:context} \\
$\vartheta = f_\xi(\mathcal{E})$
  & $(\beta,\{\alpha_j\},\gamma)$
    (Box~3.1--3.3)
  & $(\beta,\{\alpha_j\},\gamma,\lambda_s,\lambda_h)$
    \eqref{eq:ch4:p_update} \\
Primary objective $\mathcal{L}$
  & $\mathcal{L}_{\mathrm{P1}}$: expected traj.\ cost
    \eqref{eq:ch4:LP1}
  & $\mathcal{L}_{\mathrm{P2}}$: CVaR$_{0.95}(J)$
    \eqref{eq:ch4:LP2} \\
Update hierarchy $\mathcal{U}$
  & 1 loop: segment Gauss-Newton
    (Box~3.3)
  & 3 loops: segment / episode / curriculum
    (\S\ref{sec:ch4:loops}) \\
Persistent state
  & Table~\ref{tab:ch3:cache}: $\xi$, $\zeta_k$, $J_k$
  & Table~\ref{tab:ch4:state}: adds $\hat{c}_e$,
    $\Pi_{\mathrm{curr}}^e$, phase flag \\
\bottomrule
\end{tabular}}
\end{table}

% =======================================================================
\section{The Enrichment Principle: Trace First, Theorem Second}
\label{sec:ch5:principle}
% =======================================================================

\subsection{The Concrete Trace}
\label{subsec:ch5:trace}

Consider the single additive term $-\tau\nabla_\Pos H_{\mathrm{mat}}$
in the Stage~2 $p$-update.
Why is adding it not a simple extension?
Because it forces, by necessity, five downstream consequences.
We trace each one through the actual equations of
Chapters~\ref{chap:grlsnam}--\ref{chap:material}.

\medskip
\noindent\textbf{(i) New force channel.}
$H_{\mathrm{mat}} = \lambda_s\tilde{r}(\vec{c}) + \lambda_h b(\phi(\vec{c}))$
contributes two forces to the $p$-update:
$F_{\mathrm{soft}} = -\lambda_s\nabla\tilde{r}$ and
$F_{\mathrm{hard}} = -\lambda_h\nabla b(\phi)$
-- equations \eqref{eq:ch4:F_soft}--\eqref{eq:ch4:F_hard}.
These are new additive terms in the annotated boxed equation
\eqref{eq:ch4:p_update}, acting at every integration step.
They did not exist in Stage~1's $p$-update \eqref{eq:ch3:half_mom}.

\medskip
\noindent\textbf{(ii) New coefficient subspace.}
$\lambda_s$ and $\lambda_h$ are new output heads of the meta-navigator
$f_\theta$, realized as the material branch of
CoefEnergyNetMaterial (§\ref{sec:ch4:implementation}).
They occupy new rows in Table~\ref{tab:ch4:state} under
"persistent / unfrozen at P1→P2."
Setting $\lambda_s = \lambda_h = 0$ exactly recovers Stage~1
(backward compatibility, Definition~\ref{def:ch5:compat}).

\medskip
\noindent\textbf{(iii) New objective term.}
Without $\lambda_s > 0$, the force $F_{\mathrm{soft}} \equiv 0$
and the trajectory never responds to $\tilde{r}$.
The soft-risk accumulation term $w_r\sum_t\tilde{r}(\vec{c}_t)\Delta s_t$
in episode cost $J$ \eqref{eq:ch4:J} would then receive zero gradient
signal.
In the present framework, a learnable new coefficient requires an
objective term that makes its effect visible to optimization;
without one, $\lambda_s$ has no loss surface to descend.

\medskip
\noindent\textbf{(iv) New gradient pathway.}
The rollout sensitivity recursion \eqref{eq:ch4:dp_dtheta}
gains a new additive stream:
\[
  \frac{\partial\Mom_{t+1}}{\partial\theta}
  \;=\; \cdots
  \;\underbrace{-\,\tau\frac{\partial\nabla_\Pos H_{\mathrm{mat}}}{\partial\theta}}_{\text{new material stream}}.
\]
Specifically, $\partial\mathcal{L}/\partial\lambda_s =
\sum_t(\partial\mathcal{L}/\partial F_{\mathrm{soft},t})
\cdot(-\nabla\tilde{r})$
— equation \eqref{eq:ch4:dL_dlambda_s}.
This pathway did not exist in Stage~1's sensitivity recursion.
It is decoupled from the geometry and dissipation streams because
$H_{\mathrm{mat}}$ is an independent additive term.

\medskip
\noindent\textbf{(v) New update law.}
$\lambda_s$ is updated by the CVaR episode gradient (Loop~2,
§\ref{subsec:ch4:ep_loop}), \emph{not} by the segment
Gauss-Newton (Loop~1).
Why?
Because $\lambda_s$ governs tail-risk behavior of whether the policy
avoids the worst 5\% of rollouts.
The tail distribution requires a full episode batch to estimate;
a single segment residual carries no tail information.
This is a new computational job that cannot be folded into Loop~1.
Stage~1 had no episode-level loop at all (Table~\ref{tab:ch3:cache}).

\medskip
The chain is complete.
One additive line in the $p$-update costs five downstream consequences,
each traceable to a specific equation or algorithm in
Chapters~\ref{chap:grlsnam}--\ref{chap:material}.

\subsection{The Enrichment Principle}
\label{subsec:ch5:theorem}

\begin{theorem}[Learnable Energy-Enrichment Principle]
\label{thm:ch5:enrichment}
Within the class of continual Hamiltonian navigation instances
(Definition~\ref{def:ch5:instance}), suppose a new \emph{learnable}
energy term $\vartheta_{\mathrm{new}} H_{\mathrm{new}}(\Pos;\mathcal{E})$
with adaptive scalar coefficient $\vartheta_{\mathrm{new}}$ is added to
$H_\theta^{(k)}$.
Then the following consequences hold:
\begin{enumerate}[(i)]
  \item \emph{(Force channel, always.)} $F_{\mathrm{new}} = -\vartheta_{\mathrm{new}}\nabla_\Pos H_{\mathrm{new}}$
        enters the $p$-update at every integration step.
  \item \emph{(Coefficient, for learnable terms.)} The new parameter
        $\vartheta_{\mathrm{new}}$ enlarges the coefficient subspace; a
        fixed-prior term may be absorbed without a new degree of freedom.
  \item \emph{(Objective, for identifiability.)} If $\vartheta_{\mathrm{new}}$
        is to be learned from data, then $\mathcal{L}$ must contain at least
        one term whose gradient is sensitive to $F_{\mathrm{new}}$; otherwise
        $\partial\mathcal{L}/\partial\vartheta_{\mathrm{new}} = 0$ and
        $\vartheta_{\mathrm{new}}$ is not identifiable.
  \item \emph{(Sensitivity stream, always.)} The rollout sensitivity
        recursion gains a new additive stream
        $-\tau\vartheta_{\mathrm{new}}\,\partial\nabla_\Pos H_{\mathrm{new}}/\partial\theta$.
  \item \emph{(Update timescale, for distributional signals.)} If the
        gradient signal needed to identify $\vartheta_{\mathrm{new}}$ is
        only available at an episodic or distributional timescale (e.g.,
        CVaR requires a batch), a new or modified update loop in
        $\mathcal{U}$ is necessary.
\end{enumerate}
The enriched instance $\mathfrak{S}^{(k+1)}$ has the same skeleton
template \eqref{eq:ch5:template} with strictly richer
$H_\theta$, $\vartheta$, $\mathcal{L}$, and $\mathcal{U}$.
\end{theorem}

\begin{proof}[Proof sketch]
(i) Follows from Hamilton's equations: adding $\vartheta H_{\mathrm{new}}$
adds $-\vartheta\nabla_\Pos H_{\mathrm{new}}$ to $\dot{p}$.
(ii) A fixed (non-adaptive) term can be folded into $H_{\mathrm{new}}$
directly; no new coefficient is needed.
(iii) Chain rule: $\partial\mathcal{L}/\partial\vartheta_{\mathrm{new}}
= \sum_t (\partial\mathcal{L}/\partial F_t) \cdot F_{\mathrm{new},t}$,
which is zero if $\mathcal{L}$ is insensitive to $F_{\mathrm{new}}$.
(iv) Direct differentiation of the $p$-update recursion.
(v) If the signal requires aggregating over episodes (e.g., a tail
quantile), it is not available within a single segment loop, necessitating
a higher-frequency loop.
\end{proof}

\begin{corollary}[Concrete instantiation]
\label{cor:ch5:concrete}
Every consequence (i)--(v) of Theorem~\ref{thm:ch5:enrichment}
is instantiated in the Stage~1→Stage~2 transition:
(i) \eqref{eq:ch4:F_soft}--\eqref{eq:ch4:F_hard} and
\eqref{eq:ch4:p_update};
(ii) $(\lambda_s,\lambda_h)$ in Table~\ref{tab:ch4:state};
(iii) $J$ terms 3--4 in \eqref{eq:ch4:J};
(iv) \eqref{eq:ch4:dL_dlambda_s}--\eqref{eq:ch4:dL_dlambda_h};
(v) Loop~2 in §\ref{subsec:ch4:ep_loop} and its absence from
    Stage~1's Box~3.3.
\end{corollary}

\begin{proposition}[Hamiltonian risk-potential duality]
\label{prop:ch5:risk_duality}
At Stage~2 convergence, the material energy
$H_{\mathrm{mat}}(\vec{c})
 = \lambda_s^*\tilde{r}(\vec{c})
 + \lambda_h^*b_{\mathrm{sp}}(\phi(\vec{c}))$
is a risk potential in the Hamiltonian sense:
the conservative force
$F_{\mathrm{mat}} = -\nabla_{\vec{q}}H_{\mathrm{mat}}$
coincides with the negative gradient of expected soft-risk exposure
along trajectories.
Consequently, Hamiltonian dynamics minimise risk exposure not through
explicit replanning but because risk is encoded as potential energy where
the system is deflected away from high-risk regions by the same
mechanism by which a physical system evolves away from
high-potential configurations.
\end{proposition}
\begin{proof}
By Hamilton's equations,
$\dot{\Mom}=-\nabla_{\vec{q}}H$, so
$F_{\mathrm{mat}}
 = -\lambda_s^*\nabla\tilde{r}
 - \lambda_h^*\nabla b_{\mathrm{sp}}$
(\eqref{eq:ch4:F_soft}--\eqref{eq:ch4:F_hard}).
At convergence, $\partial\,\mathrm{CVaR}/\partial\lambda_s=0$ implies
$\lambda_s^*$ is the scalar making $-\lambda_s^*\nabla\tilde{r}$
the equilibrium deflection force on the optimal trajectory.
Since $J$ contains $w_r\sum_t\tilde{r}(\vec{c}_t)\Delta s_t$
(\eqref{eq:ch4:J}, term~3), minimising $\mathrm{CVaR}(J)$ over
$\lambda_s$ is equivalent to learning the slope of a risk potential
whose gradient field equals $F_{\mathrm{soft}}$.
The $\lambda_h^*$ term follows identically for the hard-hazard barrier.
\end{proof}

\begin{corollary}[Material-aware enrichment yields lower-risk path selection]
\label{cor:ch5:risk_selection}
Let $\Pi_{\mathrm{geom}}$ be the set of trajectories feasible under
Stage~1's geometric objective within the waypoint-induced routing
corridor.
At Stage~2 convergence:
\begin{enumerate}[(i)]
  \item \emph{Local preference.}
    For any $\tau_1,\tau_2\in\Pi_{\mathrm{geom}}$ with equal geometric
    cost and $\int\tilde{r}\,|_{\tau_1}>\int\tilde{r}\,|_{\tau_2}$,
    Stage~2 strictly prefers $\tau_2$.
    (Lemma~\ref{lem:ch4:risk_deflect}, applied pointwise.)
  \item \emph{Tail-risk dominance.}
    $\mathrm{CVaR}_\beta^{(\mathrm{S2})}
     \leq\mathrm{CVaR}_\beta^{(\mathrm{S1})}$
    for all $\beta\in(0,1)$, with strict inequality whenever
    $\Pi_{\mathrm{geom}}$ contains trajectories of differing risk
    exposure.
    (Proposition~\ref{prop:ch4:cvar_descent}.)
  \item \emph{Continual improvement.}
    CVaR subgradient iterates converge to a stationary point of the
    tail-risk objective at rate $O(1/\sqrt{k})$,
    reducing per-episode risk exposure toward a locally optimal level.
    (Proposition~\ref{prop:ch4:cvar_descent}.)
\end{enumerate}
Stage~2 does not globally minimise risk over all routes, it is
constrained to the feasibility corridor from the geometric scaffold.
Theorem~\ref{thm:ch5:enrichment} guarantees that the enriched structure
is necessary; this corollary guarantees it is sufficient for
lower-risk path selection within that corridor, improving continuously
over experience.
\end{corollary}

% =======================================================================
\section{The Four Coupled Enlargements}
\label{sec:ch5:four}
% =======================================================================

Table~\ref{tab:ch5:enlargements} traces each enlargement axis through
the actual equations of both stages, with the equation reference that
makes each claim checkable.

\begin{table}[t]
\centering
\caption{%
  \textbf{Four coupled enlargements from Stage~1 to Stage~2.}
  Every entry is grounded in a specific equation or section
  reference from Chapters~\ref{chap:grlsnam}--\ref{chap:material}.
}
\label{tab:ch5:enlargements}
\small
\setlength{\tabcolsep}{4pt}
\renewcommand{\arraystretch}{1.4}
\begin{tabular}{p{2.2cm}p{3.8cm}p{3.8cm}p{3.0cm}}
\toprule
\textbf{Axis} &
\textbf{Stage~1} &
\textbf{Stage~2} &
\textbf{What it enables} \\
\midrule
\textbf{A. Stored}\par\textbf{energy}
  & $H^{(1)}_{\mathrm{kin}} + H_{\mathrm{geom}}$
    \eqref{eq:ch3:ham}
  & $H^{(1)} + \lambda_s\tilde{r} + \lambda_h b(\phi)$;\;
    $R_\theta$ explicit port
    \eqref{eq:ch4:ham}, \eqref{eq:ch4:R}
  & Material-risk and hazard barrier force channels in $p$-update \\
\midrule
\textbf{B. Coefficient}\par\textbf{space}
  & $(\beta,\{\alpha_j\},\gamma)$;\;
    3 output heads in $f_\xi$
    (Box~3.1)
  & $(\beta,\{\alpha_j\},\gamma,\lambda_s,\lambda_h)$;\;
    2 new material heads
    (\S\ref{sec:ch4:implementation})
  & New gradient pathways \eqref{eq:ch4:dL_dlambda_s}--\eqref{eq:ch4:dL_dlambda_h};
    backward-compatible init \\
\midrule
\textbf{C. Objective}
  & $\mathcal{L}_{\mathrm{P1}}$: expected traj.\ cost
    \eqref{eq:ch4:LP1}
  & $\mathrm{CVaR}_{0.95}(J) + 0.3\mathcal{L}_{\mathrm{P1}}$;\;
    tail-risk on $J$ \eqref{eq:ch4:LP2};\;
    double regularization \eqref{eq:ch4:L_total}
  & Catastrophic tail penalized (Prop.~\ref{prop:ch4:cvar});
    geometry scaffold retained \\
\midrule
\textbf{D. Update}\par\textbf{hierarchy}
  & 1 loop: segment Gauss-Newton
    (Box~3.3, Table~\ref{tab:ch3:cache})
  & 3 loops: segment (local), episode (global CVaR),
    curriculum (conquest + phase)
    (\S\ref{sec:ch4:loops}, Table~\ref{tab:ch4:state})
  & Statistical coverage; passivity-gated phase transition;
    persistent conquest flags \\
\bottomrule
\end{tabular}
\end{table}

\paragraph{Reading the table.}
Row~A shows the single additive difference in $H_\theta$ that
initiates the enrichment chain.
Row~B shows the new coefficient slots that Consequence~(ii) forces.
Row~C shows the new objective terms that Consequence~(iii) forces.
Row~D shows the new update laws that Consequence~(v) forces.
The table is Table~\ref{tab:ch5:instances} written at the level of
consequence rather than component.

\paragraph{Backward compatibility.}
\begin{definition}[Backward-compatible enrichment]
\label{def:ch5:compat}
An enrichment $\mathfrak{S}^{(k)}\to\mathfrak{S}^{(k+1)}$ is
\emph{backward-compatible} if setting
$\vartheta_{\mathrm{new}} = 0$ recovers
$H^{(k+1)} = H^{(k)}$ exactly.
\end{definition}
Both stage transitions satisfy this:
$\lambda_s = \lambda_h = 0$ gives $H^{(2)} = H^{(1)}$.
The Phase~1 weights serve as a warm start for Phase~2;
Algorithm~\ref{algo:ch4:curriculum} initializes with $\lambda_s=\lambda_h=0$
frozen and only unfreezes them after the passivity criterion is met.

\paragraph{Passivity at initialization.}
\begin{proposition}[Passivity preserved at backward-compatible init]
\label{prop:ch5:passivity}
If $\mathfrak{S}^{(k)}$ satisfies
$\mathcal{L}_{\mathrm{pass}} \leq \epsilon_{\mathrm{pass}}$
at convergence, then at the backward-compatible initialization of
$\mathfrak{S}^{(k+1)}$ (with $\vartheta_{\mathrm{new}} = 0$),
$\mathcal{L}_{\mathrm{pass}}^{(k+1)} = \mathcal{L}_{\mathrm{pass}}^{(k)}
\leq \epsilon_{\mathrm{pass}}$.
\end{proposition}
\begin{proof}
At $\vartheta_{\mathrm{new}} = 0$, $H^{(k+1)} = H^{(k)}$ exactly,
so $\Delta H^{(k+1)}_t = \Delta H^{(k)}_t$ along every trajectory.
The Phase~1 curriculum only advances (Algorithm~\ref{algo:ch4:curriculum},
line~9) when $\mathcal{L}_{\mathrm{pass}} \leq \epsilon_{\mathrm{pass}}$,
so the condition is guaranteed at initialization.
\end{proof}

The spike in $\mathcal{L}_{\mathrm{pass}}$ observed early in Phase~2
training (Figure~\ref{fig:ch6:curriculum}) is due to $\lambda_s,\lambda_h$
moving away from zero, temporarily increasing the stored energy;
the curriculum delays advancement until passivity is re-established.

% =======================================================================
\section{What the Progression Achieves Computationally}
\label{sec:ch5:compute}
% =======================================================================

The clearest computational picture of the enrichment is through
the persistent-state tables.
Table~\ref{tab:ch3:cache} shows what Stage~1 maintains:
$\xi$ (fixed at inference), the segment coefficient vector $\zeta_k$,
and the smoothed Jacobian $J_k$ are the three objects, two of which are
segment-local.
Table~\ref{tab:ch4:state} shows what Stage~2 adds:
persistent conquest flags $\hat{c}_e$, a curriculum sampling
distribution $\Pi_{\mathrm{curr}}^e$, a phase flag,
and frozen-then-unfrozen material coefficients $(\lambda_s,\lambda_h)$.

The passage from Table~\ref{tab:ch3:cache} to Table~\ref{tab:ch4:state}
is the computational picture of Enlargement~D.
"More adaptation hierarchy" is not an abstract claim, it is a
concrete expansion of the maintained state across loops:
Stage~2 adds persistent conquest flags $\hat{c}_e$, a curriculum
sampling distribution $\Pi_{\mathrm{curr}}^e$, a phase flag,
and frozen-then-unfrozen material coefficients $(\lambda_s,\lambda_h)$,
each with its own update trigger and update cost.

\paragraph{Failure modes addressed.}
Table~\ref{tab:ch6:failures} in Chapter~\ref{chap:experiments}
documents the one-to-one correspondence: each added structural
component resolves a specific failure mode.
$H_{\mathrm{mat}}^{\mathrm{soft}}$ resolves hazard drift
(17→3 episodes); $R_\theta$ as explicit port resolves oscillation
(8.7→1.2 per episode); CVaR resolves catastrophic tail encounters
(0 hard-hazard entries vs.\ 7 without $H_{\mathrm{mat}}^{\mathrm{hard}}$);
the passivity curriculum resolves premature phase transition instability.
The enrichment principle predicts exactly these correspondences:
each new structural element was forced by a specific gap in the
previous stage's machinery.

% =======================================================================
\section{Current Limits and the Next Rung}
\label{sec:ch5:limits}
% =======================================================================

Each limitation of Stage~2 is expressible as a missing component
in Table~\ref{tab:ch4:state}.

\paragraph{Oracle material access.}
$(\lambda_s,\lambda_h)$ are trained on oracle $\tilde{r},\phi$.
Stage~3 replaces these with learned belief $\hat\rho$ from RGB-IR,
adding a new entry to the persistent state (perception backbone
weights $\omega$) and a new gradient pathway
$\partial\mathcal{L}/\partial\omega$ through the perception model.
The enrichment chain for Stage~3 is identical in structure to the
Stage~1→2 chain:
$H_{\mathrm{perc}} \to F_{\mathrm{perc}} \to \omega \to
\mathcal{L}_{\mathrm{perc}} \to \nabla_\omega\mathcal{L} \to$
belief update loop.

\paragraph{Simplified observer.}
The context update $y_{t+1} = \Psi(y_t,\vec{q}_{t+1})$
carries no belief state.
A richer observer would add a learned $b_t$ to the persistent state
and a filtering update law to $\mathcal{U}$ -- a fourth loop.

\paragraph{Isotropic metric.}
$\Mass$ is a fixed scalar.
A learned $\Mass(\vec{q};\xi)$ would add position-dependent entries
to the kinetic term and a new gradient stream
$\partial H_{\mathrm{kin}}/\partial\xi$ through the $p$-update.

\paragraph{Empirical coverage.}
The conquest flags $\hat{c}_e$ are a finite-sample proxy.
The theoretical open problem (§\ref{sec:ch7:theory} T3) is bounding
$\mathrm{vol}(\mathcal{B}(\theta))$ analytically -- replacing the
finite-sample flag with a Riemannian volume estimate.

% =======================================================================
\section{Chapter Summary}
\label{sec:ch5:summary}
% =======================================================================

The synthesis of this chapter can be read directly off the costate
equation \eqref{eq:ch5:side_by_side}.
Stage~2 adds one term, $-\tau\nabla_\Pos H_{\mathrm{mat}}$,
to the Stage~1 $p$-update.
That single term is not a simple extension because it forces,
by necessity:
a new force channel in $F_{\mathrm{tot}}$ (i),
a new coefficient in $f_\xi$'s output (ii),
a new term in episode cost $J$ (iii),
a new gradient pathway in the sensitivity recursion (iv),
and a new computational loop in $\mathcal{U}$ (v).

Theorem~\ref{thm:ch5:enrichment} states this abstractly.
Corollary~\ref{cor:ch5:concrete} verifies it by pointing to specific
equations.
Table~\ref{tab:ch5:enlargements} traces all four enlargement axes
with equation cross-references.
The passage from Table~\ref{tab:ch3:cache} to Table~\ref{tab:ch4:state}
gives the concrete computational picture of what "richer hierarchy"
means.

The enrichment principle is \emph{generative}: it predicts the
structure of Stage~3 (perception backbone, belief update loop) in
exactly the same way it characterizes Stage~2.
The ladder's organizing rule: one new energy term costs five
downstream consequences is what makes the framework a research
program rather than a pair of navigation systems.

\begin{figure}[t]
\centering
% Pre-define named colors to avoid nested xcolor mixing in parameterized TikZ styles
\colorlet{stagegray}{gray!60}
\begin{tikzpicture}[
  font=\small\sffamily,
  % Stage box style
  stage/.style={
    rectangle, rounded corners=6pt,
    draw=#1!70!black, fill=#1!9,
    line width=1.4pt,
    minimum width=13.5cm, minimum height=2.0cm,
    align=left, inner sep=10pt
  },
  stageno/.style={
    circle,
    draw=#1!80!black, fill=#1!25,
    line width=1.2pt,
    minimum size=26pt,
    font=\bfseries\normalsize\sffamily
  },
  % Chain node style
  chain/.style={
    rectangle, rounded corners=3pt,
    draw=#1!60!black, fill=#1!12,
    line width=0.8pt,
    minimum width=1.9cm, minimum height=0.7cm,
    align=center, inner sep=3pt,
    font=\scriptsize\sffamily
  },
  % Arrow styles
  harrow/.style={-{Stealth[length=5pt]}, line width=1pt, color=gray!60},
  varrow/.style={-{Stealth[length=7pt]}, line width=1.4pt, color=#1},
  addlbl/.style={font=\scriptsize\sffamily\itshape, text=gray!60!black},
  % Dashed future
  futurebox/.style={
    rectangle, rounded corners=6pt,
    draw=gray!40, fill=gray!5,
    line width=1pt, dashed,
    minimum width=13.5cm, minimum height=1.6cm,
    align=left, inner sep=10pt
  },
]

% ── Stage 0: Geometry-only (Ch 3) ─────────────────────────────────────
\node[stage=teal] (s0) at (0,0) {%
  \hspace{30pt}%
  \begin{minipage}{12.2cm}
    \textbf{Stage 1 — Geometry-Only (Chapter~\ref{chap:grlsnam})}\\[3pt]
    $H^{(1)} = H_{\mathrm{kin}} + \beta D_\mathcal{G}
               + \sum_j\alpha_j b(d_j)$,\quad
    $R^{(1)}: \gamma\Mass^{-1}\Mom$ (implicit)\\[2pt]
    $\mathcal{L}^{(1)} = \mathcal{L}_{\mathrm{geom}}$,\quad
    $\mathcal{U}^{(1)}$: 1-scale online $\alpha_j$ update\\[2pt]
    \textit{Solves:} geometry-aware navigation under local sensing.
    \textit{Cannot:} reason about latent material risk.
  \end{minipage}
};
\node[stageno=teal] at ($(s0.west)+(18pt,0)$) {\textcolor{teal!80!black}{1}};

% ── Transition annotation ────────────────────────────────────────────
\node[addlbl] at ($(s0.south)!0.5!(s1.north)+(0,-0.18)$) {};

% % ── Enrichment chain (floating between stages) ─────────────────────
% \node[chain=orange] (cH)  at (-5.2,-2.05) {$+H_{\mathrm{mat}}$};
% \node[chain=orange] (cF)  at (-2.7,-2.05) {$\to F_{\mathrm{mat}}$};
% \node[chain=orange] (cV)  at (-0.2,-2.05) {$\to \vartheta_{\mathrm{new}}$};
% \node[chain=orange] (cL)  at ( 2.3,-2.05) {$\to \mathcal{L}_{\mathrm{new}}$};
% \node[chain=orange] (cU)  at ( 4.8,-2.05) {$\to \text{update}$};
% \draw[harrow] (cH.east) -- (cF.west);
% \draw[harrow] (cF.east) -- (cV.west);
% \draw[harrow] (cV.east) -- (cL.west);
% \draw[harrow] (cL.east) -- (cU.west);
% \node[addlbl, above=2pt of cH] {\textbf{Enrichment chain} (Theorem~\ref{thm:ch5:enrichment})};

% ── Stage 1: Material-aware (Ch 4) ──────────────────────────────────
\node[stage=orange, below=1.5cm of s0] (s1) {%
  \hspace{30pt}%
  \begin{minipage}{12.2cm}
    \textbf{Stage 2 — Material-Aware (Chapter~\ref{chap:material})}\\[3pt]
    $H^{(2)} = H^{(1)} + \lambda_s\tilde{r}(\vec{c})
               + \lambda_h b(\phi(\vec{c}))$,\quad
    $R^{(2)}: \gamma\Mass^{-1}\Mom$ (explicit port)\\[2pt]
    $\mathcal{L}^{(2)} = \mathrm{CVaR}_{0.95}(J) +
               \mathcal{R}_H + \mathcal{R}_{\nabla H}$,\quad
    $\mathcal{U}^{(2)}$: segment $\to$ episode $\to$ curriculum\\[2pt]
    \textit{Solves:} risk-sensitive navigation under oracle material.
    \textit{Cannot:} infer material from raw perception.
  \end{minipage}
};
\node[stageno=orange] at ($(s1.west)+(16pt,0)$)
  {\textcolor{orange!80!black}{2}};

\draw[varrow=orange!70!black] (s0.south) -- (s1.north);

% % ── Enrichment chain (between s1 and s2_future) ─────────────────────
% \node[chain=gray] (dH) at (-5.2,-4.1) {\small$+H_{\mathrm{perc}}$};
% \node[chain=gray] (dF) at (-2.7,-4.1) {\small$\to F_{\mathrm{perc}}$};
% \node[chain=gray] (dV) at (-0.2,-4.1) {\small$\to \vartheta_{\mathrm{perc}}$};
% \node[chain=gray] (dL) at ( 2.3,-4.1) {\small$\to \mathcal{L}_{\mathrm{perc}}$};
% \node[chain=gray] (dU) at ( 4.8,-4.1) {\small$\to \text{update}$};
% \draw[harrow] (dH.east) -- (dF.west);
% \draw[harrow] (dF.east) -- (dV.west);
% \draw[harrow] (dV.east) -- (dL.west);
% \draw[harrow] (dL.east) -- (dU.west);
% \node[addlbl, above=2pt of dH, text=gray!60]
%   {\textit{next enrichment (future)}};

% ── Stage 2: Future / deferred ──────────────────────────────────────
\node[futurebox, below=1.5cm of s1] (s2) {%
  \hspace{30pt}%
  \begin{minipage}{12.2cm}
    \textbf{Stage 3 — Perception-Led Material Inference (Future Work)}\\[3pt]
    $H^{(3)} = H^{(2)} + H_{\mathrm{perc}}(\vec{c}; \hat\rho)$,\quad
    learned material belief $\hat\rho$ from RGB-IR\\[2pt]
    $\mathcal{L}^{(3)}$ adds perception loss, uncertainty calibration;
    $\mathcal{U}^{(3)}$ adds belief update loop\\[2pt]
    \textit{Would solve:} end-to-end perception-to-navigation under
    latent material uncertainty.
  \end{minipage}
};
\node[stageno=stagegray] at ($(s2.west)+(16pt,0)$)
  {\textcolor{gray!50}{\textbf{3}}};
\draw[varrow=gray!50, dashed] (s1.south) -- (s2.north);

% ── Right-side axis label ──────────────────────────────────────────
\draw[{Stealth[length=6pt]}-{Stealth[length=6pt]},
      line width=1pt, color=gray!50]
  ($(s0.north east)+(0.5cm,0)$) -- ($(s2.south east)+(0.5cm,0)$)
  node[midway, right=6pt, font=\small\sffamily\bfseries,
       text=gray!60, rotate=-90, anchor=south] {enrichment axis};

\end{tikzpicture}
\caption{%
  \textbf{The enrichment ladder.}
  Each transition arrow is labeled with the five-link chain
  $(H_{\mathrm{new}} \to F \to \vartheta \to \mathcal{L} \to \nabla\mathcal{L} \to \text{loop})$.
  Stage~1 maintains segment-local state only;
  Stage~2 adds persistent curriculum state across episodes;
  Stage~3 (dashed, future) adds a perception backbone to the
  maintained state.
  The same enrichment principle governs each transition.
}
\label{fig:ch5:enrichment_ladder}
\end{figure}